\documentclass[11pt,a4paper]{article}

% ---- Packages ----
\usepackage[margin=2.5cm]{geometry}
\usepackage[utf8]{inputenc}
\usepackage[T1]{fontenc}
\usepackage{lmodern}
\usepackage{amsmath,amssymb,amsthm}
\usepackage{graphicx}
\usepackage[dvipsnames]{xcolor}
\usepackage{booktabs}
\usepackage{tabularx}
\usepackage{multirow}
\usepackage{enumitem}
\usepackage{caption}
\usepackage[numbers,sort&compress]{natbib}
\usepackage{lineno}
\usepackage{fancyhdr}
\usepackage[hidelinks,colorlinks=true,linkcolor=blue!60!black,citecolor=blue!60!black,urlcolor=blue!60!black]{hyperref}
\usepackage{microtype}
\usepackage{mdframed}

% ---- Theorem environments ----
\newtheorem{theorem}{Theorem}

% ---- Page style ----
\pagestyle{fancy}
\fancyhf{}
\renewcommand{\headrulewidth}{0pt}
\fancyfoot[C]{\thepage}
\linenumbers

% ---- Paragraph style ----
\setlength{\parindent}{0pt}
\setlength{\parskip}{6pt}

% ---- Custom commands ----
\newcommand{\norm}[1]{\left\|#1\right\|}
\newcommand{\eps}{\varepsilon}

\begin{document}

\begin{center}
{\LARGE\bfseries Supplementary Information}\\[8pt]
{\large Designing Any Imaging System from Natural Language:\\
Agent-Constrained Composition over a Finite Primitive Basis}\\[12pt]
{\normalsize Chengshuai Yang}\\[6pt]
\today
\end{center}

\vspace{12pt}

% ========================================================================
\section{S1. Foundational Theorems from the Companion Paper}
% ========================================================================

The companion paper~\cite{paperII} establishes two foundational results that this paper builds upon.  To make the present work self-contained, we reproduce the full theorem statements, proof sketches, and key lemmas below.

\subsection{S1.1. Finite Primitive Basis (FPB) Theorem}

\begin{theorem}[Finite Primitive Basis]
\label{thm:fpb_supp}
For every imaging forward model $H$ in the designable class $\mathcal{C}_{\mathrm{tier}}$, there exists a well-formed typed DAG $G = (V, E, \tau)$ whose node types are drawn from the canonical primitive library $\mathcal{B} = \{P, M, \Pi, F, C, \Sigma, D, S, W, R, \Lambda\}$ (11 elements) such that $G$ is an $\eps$-approximate representation of $H$ with $\eps = 0.01$.
\end{theorem}

\textbf{Proof sketch.}
The proof is constructive.  Given any forward model $H = H_K \circ H_{K-1} \circ \cdots \circ H_1$, we decompose it into a sequence of physics stages and show each stage is representable by one or more FPB primitives via six \emph{Primitive Realization Lemmas}:

\textbf{Lemma 1 (Propagation Realization).}
Free-space wave propagation $H_k$ is represented by the Propagate primitive $P(d, \lambda)$ via the angular spectrum representation: $P = \mathcal{F}^{-1}[T_{\mathrm{exact}} \cdot \mathcal{F}[\cdot]]$.  Error from evanescent wave truncation: $\eps_{\mathrm{evan}} = e^{-2\pi d/\lambda}$, negligible for $d \gg \lambda$.

\textbf{Lemma 2 (Elastic Interaction Realization).}
Element-wise transmission multiplication is \emph{exactly} represented by the Modulate primitive $M(\mathbf{m})$: $\norm{H_k - M(\mathbf{m})} = 0$.

\textbf{Lemma 3 (Scattering Realization).}
Under the first Born approximation, the scattered field is a linear functional representable by the Scatter primitive $R$.  Single-scattering representation is exact; $L$-th order multiple scattering error: $\eps_{\mathrm{scat}}^{(L)} \leq (\norm{V}/k)^{L+1}$ (geometric convergence).

\textbf{Lemma 4 (Encoding-Projection Realization).}
Line-integral projection (CT) is exactly the Radon transform $\Pi(\theta)$; Fourier encoding (MRI) is exactly $F(\mathbf{k})$.  Both are exact: $\eps = 0$.

\textbf{Lemma 5 (Detection-Readout Realization).}
The detection chain decomposes into $\leq$5 sequential operations, each representable by a primitive.  Detector PSF error: $\eps_{\mathrm{PSF}} \leq \ell_c^2/p^2$ (crosstalk length $\ell_c$, pixel pitch $p$).  Total: $\eps_{\mathrm{det}} < 0.01$ for all 39 validated modalities.

\textbf{Lemma 6 (Nonlinear-Stage Realization).}
All pointwise physics nonlinearities belong to 5 structural categories (exponential, logarithmic, phase wrapping, polynomial, saturation), each representable by the Transform primitive $\Lambda(f, \boldsymbol{\theta})$.

\textbf{Error composition bound.}
Let $\gamma_k = \norm{H_k}$ (linear) or $\mathrm{Lip}(H_k)$ (nonlinear).  By Lipschitz composition:
\[
\sup_{\norm{\mathbf{x}} \leq 1} \norm{H(\mathbf{x}) - H_{\mathrm{dag}}(\mathbf{x})} \leq \sum_{k=1}^{K} \eps_k \prod_{j=k+1}^{K} \gamma_j \leq K \cdot \max_k(\eps_k) \cdot B^{K-1}
\]
With empirical maximum $K{=}5$ and $B{=}4$: requires $\max_k(\eps_k) \leq 3.9 \times 10^{-5}$, satisfied by all 6 lemmas.  Per-modality tightening exploits the fact that most $\eps_k = 0$ (elastic, encoding stages) and most primitives are norm-preserving.

\textbf{Necessity.}
All 11 primitives are necessary: for each primitive $B_i$, there exists a witness modality whose representation error exceeds 0.01 when $B_i$ is removed.  Witnesses: $P$ $\leftarrow$ ptychography, $M$ $\leftarrow$ CASSI, $\Pi$ $\leftarrow$ CT, $F$ $\leftarrow$ MRI, $C$ $\leftarrow$ lensless, $\Sigma$ $\leftarrow$ SPC, $D$ $\leftarrow$ all, $S$ $\leftarrow$ MRI, $W$ $\leftarrow$ CASSI, $R$ $\leftarrow$ Compton ($\eps = 0.34$ without), $\Lambda$ $\leftarrow$ polychromatic CT.

\textbf{Complexity bound.}
Node count: $|V| \leq 20$ (empirical max: 5).  Depth: $\mathrm{depth}(G) \leq 10$ (empirical max: 5).

\subsection{S1.2. Triad Decomposition Theorem}

\begin{theorem}[Triad Decomposition]
\label{thm:triad_supp}
Let $H \in \mathcal{C}_{\mathrm{tier}}$ with measurement $\mathbf{y} = H\mathbf{x} + \mathbf{n}$ and reconstruction using nominal operator $H_{\mathrm{nom}}$.  The reconstruction MSE decomposes as:
\[
\mathrm{MSE} \leq \mathrm{MSE}^{(G_1)} + \mathrm{MSE}^{(G_2)} + \mathrm{MSE}^{(G_3)}
\]
where $\mathrm{MSE}^{(G_1)}$ is the null-space loss (information-theoretic floor), $\mathrm{MSE}^{(G_2)}$ is the noise-floor term (carrier budget), and $\mathrm{MSE}^{(G_3)}$ is the mismatch-induced term (model fidelity).  This decomposition is diagnostically complete: every reconstruction failure is attributable to at least one gate.
\end{theorem}

\textbf{Gate 1 (Recoverability / Compression Bound).}
For any estimator $\hat{\mathbf{x}}(\mathbf{y})$, the minimum achievable MSE satisfies:
\[
\mathrm{MSE}_{\min}^{(G_1)} \geq \frac{1}{n}\sum_{i=1}^{n-m} \sigma_i^2(\mathbf{x})
\]
where $\sigma_i^2(\mathbf{x})$ is the variance of $\mathbf{x}$ along the $i$-th null-space direction of $H$, $m$ is the number of measurements, and $n$ is the signal dimension.  The corresponding PSNR ceiling is:
\[
\mathrm{PSNR}_{\max}^{(G_1)} = 10\log_{10}\!\left(\frac{\norm{\mathbf{x}}_\infty^2}{\mathrm{MSE}_{\min}^{(G_1)}}\right)
\]
No algorithm can exceed this ceiling without additional priors.

\textbf{Gate 2 (Carrier Budget / Noise Bound).}
Under matched forward model and Gaussian noise $\mathbf{n} \sim \mathcal{N}(0, \sigma_n^2 \mathbf{I})$:
\[
\mathrm{PSNR}_{\max}^{(G_2)} = 10\log_{10}(\mathrm{SNR}) + C_{\mathcal{M}}, \quad C_{\mathcal{M}} = 10\log_{10}\!\left(\frac{n\norm{\mathbf{x}}_\infty^2}{\norm{\mathbf{x}}^2} \cdot \kappa^{-2}(H)\right)
\]
where $\kappa(H)$ is the condition number of $H$ restricted to its column space.  No algorithm can exceed this ceiling regardless of computational budget.

\textbf{Gate 3 (Operator Mismatch / Calibration Sensitivity).}
For small perturbation $\delta\boldsymbol{\theta}$ around true parameters $\boldsymbol{\theta}^*$:
\[
\Delta\mathrm{PSNR} \approx -\frac{10}{\ln 10} \frac{\delta\boldsymbol{\theta}^\top \mathbf{J}^\top \mathbf{J}\, \delta\boldsymbol{\theta}}{\mathrm{MSE}_0}
\]
where $\mathbf{J} = \partial(H_{\boldsymbol{\theta}}\mathbf{x})/\partial\boldsymbol{\theta}\big|_{\boldsymbol{\theta}^*}$ is the parameter Jacobian and $\mathrm{MSE}_0$ is the noise-only MSE at $\boldsymbol{\theta}^*$.

\textbf{Gate 3 dominance condition.}
For instruments operating above Gates~1 and~2 floors, Gate~3 becomes the binding constraint when $\norm{\delta\boldsymbol{\theta}}_{\mathbf{J}^\top\mathbf{J}} > (\sigma_n / \norm{\mathbf{x}}_\infty) \cdot \kappa(H)$.  Empirically, Gate~3 dominates in all 12 validated modalities across 5 carrier families: CASSI shows 13.98\,dB degradation from sub-pixel mask shift; ptychography shows 42.06\,dB sensitivity to probe mismatch.

\textbf{Recovery ratio.}
$\rho = (\mathrm{PSNR}_{\mathrm{achieved}} - \mathrm{PSNR}_{\mathrm{mismatch}}) / (\mathrm{PSNR}_{\mathrm{ideal}} - \mathrm{PSNR}_{\mathrm{mismatch}})$, bounded in $[0, 1]$ under convex reconstruction loss.  Across validated modalities: $\rho \in [0.4, 0.9]$, mean 0.85.

% ========================================================================
\section{S2. Novel System Design Specifications}
% ========================================================================

The 10 novel system designs compose FPB primitives into chains not found in prior literature.  Each is specified via the 8-field \texttt{spec.md} format.  The notation $\Phi_z$ denotes depth-parameterised convolution: the $C$ primitive applied with a depth-dependent point spread function $\Phi(z)$, generated by wave propagation through a random phase diffuser with defocus modulation.

\subsection{S2.1. 2D Systems}

\textbf{Lensless imaging} ($C \to D$, compression 1:1).
A thin random phase mask creates a caustic PSF with broad spatial-frequency support.  Reconstruction inverts the convolution via ADMM+TV.  PSNR: 43.7\,dB, SSIM: 0.984.

\begin{verbatim}
modality: lensless_imaging
carrier: photon
geometry: single_shot, 128x128
object: 128x128 2D image
forward_model: Convolve(C, psf=phase_mask) -> Detect(D)
noise: Poisson I_0=5000 + Gaussian sigma=3
target: PSNR >= 40 dB
system_elements:
  source: broadband LED
  optics: random phase mask (feature_scale=2.5)
  detector: CMOS 128x128
\end{verbatim}

\textbf{SIM} ($M \to C \to D$, super-resolution).
Structured illumination modulates the sample with sinusoidal patterns, shifting high-frequency content into the passband.  Three pattern orientations at three phases yield 9 measurements for 2$\times$ resolution enhancement.  PSNR: 27.7\,dB.

\subsection{S2.2. 3D Systems (8:1 compression)}

\textbf{3D Lensless} ($\Phi_z \to \Sigma \to D$).
A random phase diffuser creates depth-dependent PSFs through defocus-modulated wave propagation: $\text{PSF}(z) = |\mathcal{F}\{\exp(i\phi_{\text{diffuser}} + i \cdot \text{defocus}(z) \cdot r^2)\}|^2$.  Each depth plane produces a unique speckle pattern; the 2D sensor measurement is the sum of all convolved depth planes.  This is the $C$ primitive instantiated per depth plane.  PSNR: 20.3\,dB at 8:1 compression.

\begin{verbatim}
modality: 3d_lensless
carrier: photon
geometry: single_shot, 128x128, n_depths=8
object: 128x128x8 3D volume (sparse depth planes)
forward_model: Convolve_z(C, psf=diffuser(z)) -> Sum(Sigma) -> Detect(D)
noise: Poisson I_0=5000 + Gaussian sigma=3
target: PSNR >= 15 dB
system_elements:
  source: broadband LED
  optics: random phase diffuser (feature_scale=1.0, defocus_max=40)
  detector: CMOS 128x128
\end{verbatim}

\textbf{Temporal-coded lensless} ($M \to C \to \Sigma \to D$).
A binary DMD mask modulates each of 8 video frames before diffuser convolution and temporal integration during a single exposure.  The mask provides per-pixel temporal diversity.  PSNR: 31.6\,dB.

\textbf{Spectral lensless} ($M \to W \to C \to \Sigma \to D$).
A coded mask followed by prism dispersion separates 8 spectral bands before diffuser convolution and spectral integration.  PSNR: 36.3\,dB.

\subsection{S2.3. 4D Systems (16:1 compression)}

\textbf{4D Spectral-Depth} ($M \to W_\lambda \to \Phi_z \to \Sigma \to D$).
Combines coded mask, spectral dispersion, and diffuser depth encoding to recover a $(z, \lambda)$ datacube ($4 \times 4 = 16$ planes) from a single 2D shot.  PSNR: 23.9\,dB.

\textbf{4D Temporal DMD} ($M \to \Phi_z \to \Sigma \to D$).
DMD binary masks provide temporal coding; diffuser PSFs encode depth.  Recovers $(z, t)$ datacube at 16:1.  PSNR: 25.4\,dB.

\textbf{4D Temporal Streak} ($M \to W_t \to \Phi_z \to \Sigma \to D$).
Passive temporal dispersion (streak camera) replaces active DMD coding.  A fixed mask provides spatial diversity.  PSNR: 25.3\,dB.

\subsection{S2.4. 5D Systems (64:1 compression)}

\textbf{5D Full DMD} ($M \to W_\lambda \to \Phi_z \to \Sigma \to D$).
Full 5D recovery of $(x, y, z, \lambda, t)$ from a single shot at $4 \times 4 \times 4 = 64$:1 compression using coded mask + spectral dispersion + diffuser depth encoding.  PSNR: 29.9\,dB.

\textbf{5D Full Streak} ($M \to W_\lambda \to W_t \to \Phi_z \to \Sigma \to D$).
Replaces DMD with passive streak temporal dispersion.  The longest chain in the FPB framework (5 operators before detection).  PSNR: 29.9\,dB.

\subsection{S2.5. Design Patterns}

Three patterns emerge from the novel designs:

\begin{enumerate}[nosep]
\item \textbf{Active modulation (DMD) $>$ passive dispersion}: Binary masks provide per-pixel measurement diversity that continuous dispersion cannot match.  DMD-based temporal coding (31.6\,dB) outperforms streak-based (25.3\,dB) at equal compression.
\item \textbf{Algorithm choice is critical for compressive systems}: FISTA+TV or R-L dominate across novel modalities.  Inter-method PSNR CoV reaches 42--48\% for compressive systems vs.\ $<$6\% for well-conditioned ones.
\item \textbf{Graceful compression degradation}: 1:1 $\to$ 43.7\,dB; 8:1 $\to$ 20--36\,dB; 16:1 $\to$ 24--25\,dB; 64:1 $\to$ 29--30\,dB.
\end{enumerate}

% ========================================================================
\section{S3. Maskless Control Group}
% ========================================================================

To validate that spatial modulation ($M$) is a prerequisite for compressive dispersion-based imaging, we tested 9 configurations that combine depth encoding ($\Phi_z$) with spectral ($W_\lambda$) and/or temporal ($W_t$) dispersion but \emph{without} a coded mask.  All 9 configurations compile through the structural compiler (Gates~1--5) but are flagged by the Judge Agent (Gate~1, recoverability) as severely ill-conditioned:

\begin{itemize}[nosep]
\item Condition number $> 10^6$ for all maskless configurations
\item Reconstruction PSNR: 6--9\,dB (near noise floor)
\item Adding a binary coded mask ($M$) improves PSNR by 5--15\,dB at equal compression
\end{itemize}

This confirms a fundamental principle: dispersion operators ($W$) provide only global spectral/temporal diversity, while coded masks ($M$) provide per-pixel spatial diversity.  Both are necessary for well-conditioned compressive recovery.

% ========================================================================
\section{S4. Encoder Comparison for Depth Encoding}
% ========================================================================

Three physically realisable depth-encoding mechanisms were compared across 5 modulation combinations ($3 \times 5 = 15$ configurations total):

\textbf{Diffuser} ($\Phi_z^{\text{diff}}$): Random phase plate with defocus-dependent PSFs.  $\text{PSF}(z) = |\mathcal{F}\{\exp(i\phi_{\text{random}} + i \cdot \text{defocus}(z) \cdot r^2)\}|^2$.

\textbf{Multi-lens} ($L_z$): Microlens array producing depth-dependent parallax sub-pixel shifts.

\textbf{Meta-surface} ($\Psi_z$): Engineered phase pattern (cubic + trefoil + astigmatism aberrations) with defocus modulation.

\begin{table}[htbp]
\centering
\caption{Depth encoder comparison: best PSNR (dB) across 3 encoders $\times$ 5 modulation combinations.  $N_z{=}8$ depth planes, $N_\lambda{=}4$ spectral bands, $N_t{=}4$ temporal frames, image size $128{\times}128$.}
\small
\begin{tabular}{@{}lccccc@{}}
\toprule
\textbf{Encoder} & \textbf{3D only} & \textbf{+Mask} & \textbf{+M+$W_\lambda$} & \textbf{+M+$W_t$} & \textbf{+M+$W_\lambda$+$W_t$} \\
& (8:1) & (8:1) & (32:1) & (32:1) & (128:1) \\
\midrule
Diffuser & 17.0 & 20.3 & 24.6 & 26.9 & 38.1 \\
Multi-lens & 16.6 & \textbf{44.1} & \textbf{41.0} & \textbf{36.1} & \textbf{38.7} \\
Meta-surface & 16.6 & 16.2 & 24.6 & 25.8 & 38.0 \\
\bottomrule
\end{tabular}
\end{table}

\textbf{Key findings:}
\begin{itemize}[nosep]
\item Without coded mask: all three encoders perform similarly (16.6--17.0\,dB), limited by the ill-conditioning of depth-only recovery at 8:1 compression.
\item Multi-lens + mask dominates: parallax shifts are simple, well-conditioned operators (sub-pixel translations).  Combined with a binary mask's per-pixel diversity, the composite forward operator is nearly diagonal, yielding 44.1\,dB.
\item At maximal compression (128:1, 5D), all encoders converge to ${\sim}$38\,dB because the mask and dispersion operators dominate the measurement diversity.
\end{itemize}

\textbf{PSF diversity analysis} (mean inter-depth cross-correlation):
\begin{itemize}[nosep]
\item Diffuser: 0.29 (moderate diversity, structured speckle)
\item Multi-lens: 0.17 (lowest correlation, shift-based diversity)
\item Meta-surface: 0.46 (highest correlation, engineered but similar patterns)
\end{itemize}

% ========================================================================
\section{S5. Reconstruction Details}
% ========================================================================

\subsection{S5.1. Algorithms}

Five reconstruction algorithms were evaluated for each novel modality:

\textbf{Wiener deconvolution}: Frequency-domain least-squares with noise regularisation.  Fastest ($<$0.1\,s) but lowest quality for compressive systems.

\textbf{GAP-TV}: Generalised alternating projection with total variation regularisation~\cite{meng2020gap}.  Alternates between data projection and TV denoising.

\textbf{FISTA+TV}: Fast iterative shrinkage-thresholding with TV proximal operator~\cite{beck2009fista}.  Nesterov-accelerated gradient descent with $O(1/k^2)$ convergence.

\textbf{ADMM+TV}: Alternating direction method of multipliers with TV splitting~\cite{boyd2011admm}.  Variable splitting separates data fidelity from regularisation.

\textbf{Richardson-Lucy (R-L)}: Maximum-likelihood deconvolution for Poisson noise.  Multiplicative updates preserve non-negativity.

\subsection{S5.2. Forward Model Implementation}

All forward models are implemented as linear operators with explicit adjoint computation.  The measurement for a multi-dimensional modality with depth encoding, mask, and dispersion is:
\[
y = \sum_{z,\lambda,t} D\left[ \text{shift}_{(\Delta x_\lambda, \Delta y_t)} \left( M \odot \text{conv}(x_{z,\lambda,t}, \text{PSF}(z)) \right) \right] + n
\]
where $D$ is detection, $M$ is the binary mask, $\text{conv}$ denotes convolution with depth-dependent PSF, $\text{shift}$ implements spectral/temporal dispersion as lateral shifts, and $n$ is Poisson + Gaussian noise.

% ========================================================================
\section{S6. LLM Implementation Details}
% ========================================================================

\subsection{S6.1. Model and Prompting}

All three agents (Plan, Judge, Execute) use Claude Sonnet~4 (Anthropic, \texttt{claude-sonnet-4-20250514}) as the backbone LLM.  The system prompt provides each agent with:
\begin{itemize}[nosep]
\item The 11 FPB primitives with typed signatures and physical semantics
\item The 173-modality registry (modality name, carrier, canonical chain, typical parameters)
\item The \texttt{spec.md} schema with field descriptions and 3 worked examples (CT, MRI, CASSI)
\item Agent-specific role instructions (Plan: generate specification; Judge: validate via Triad gates; Execute: select algorithm and predict quality)
\end{itemize}

Temperature is set to 0 for reproducibility.  The Plan Agent uses chain-of-thought prompting: it first identifies the carrier family, then selects primitives left-to-right along the physical signal path, then fills in parameters from the registry or physical constraints.  The Judge Agent applies deterministic gate checks (no LLM inference for Gates~1--4); only Gate~5 (cost estimation) and Gate~6 (FPB fidelity) use LLM reasoning.

\subsection{S6.2. Reproducibility and Robustness}

Three \texttt{spec.md} files (CT, MRI, CASSI) were independently generated 5 times each; all 15 specifications were identical, yielding PSNR within $\pm 0.3$\,dB.  The 4\% misspecification rate (7/173 modalities) occurs exclusively for modalities with ambiguous natural-language descriptions (e.g., ``4D-STEM'' could refer to scanning or convergent-beam modes).

\textbf{LLM dependency is a limitation.}  The agents have been tested only with Claude Sonnet~4.  Because 4 of 6 Judge gates are deterministic (no LLM inference), and the Execute Agent's algorithm selection is from a fixed catalog, the LLM dependency is concentrated in the Plan Agent's specification generation.  Cross-LLM robustness testing is future work; we expect that any LLM capable of structured JSON output and physics reasoning would produce comparable results for well-documented modalities, but novel or ambiguous modalities may show higher variance.

% ========================================================================
\section*{References}
% ========================================================================

\begingroup
\renewcommand{\section}[2]{}
\begin{thebibliography}{10}

\bibitem{paperII}
Yang, C.
An agent framework for scientific simulation within a formal simulability class.
\textit{Preprint} (2026).
Available at \url{https://github.com/integritynoble/Physics_World_Model}.

\bibitem{meng2020gap}
Meng, Z., Ma, J., \& Yuan, X.
End-to-end low cost compressive spectral imaging with spatial-spectral self-attention.
\textit{Proc.\ ECCV} (2020).

\bibitem{beck2009fista}
Beck, A. \& Teboulle, M.
A fast iterative shrinkage-thresholding algorithm for linear inverse problems.
\textit{SIAM J.\ Imaging Sci.} \textbf{2}, 183--202 (2009).

\bibitem{boyd2011admm}
Boyd, S., Parikh, N., Chu, E., Peleato, B. \& Eckstein, J.
Distributed optimization and statistical learning via the alternating direction method of multipliers.
\textit{Found.\ Trends Mach.\ Learn.} \textbf{3}, 1--122 (2011).

\end{thebibliography}
\endgroup

\end{document}
